\documentclass[11pt]{article}
\usepackage[T1]{fontenc}
\usepackage{lmodern,amsmath,amssymb,amsthm,geometry,hyperref}
\geometry{margin=1in}
\hypersetup{hidelinks,pdfauthor={Lluis Eriksson},pdftitle={A metastable Hamming shell with exactly matched stationary marginals}}
\newtheorem{theorem}{Theorem}[section]
\newtheorem{lemma}[theorem]{Lemma}
\newtheorem{corollary}[theorem]{Corollary}
\newcommand{\E}{\mathbb E}
\newcommand{\one}{\mathbf 1}
\newcommand{\norm}[1]{\left\|#1\right\|}
\newcommand{\ip}[2]{\langle #1,#2\rangle}
\title{A metastable Hamming shell with exactly matched stationary marginals}
\author{Lluis Eriksson}
\date{18 November 2026}
\begin{document}
\maketitle
\begin{abstract}
We construct finite reversible Doob chains for Boolean noise of correlation one half whose spectral gaps tend to zero exponentially. Their normalized stationary densities have any prescribed fixed floor below one, converge exponentially to one in cubic mean, and have exactly uniform marginals through a positive fraction of the coordinates. A single Hamming sphere of exponentially small stationary mass traps the chain for exponentially many steps. The construction adjusts the potential on the shells typically reached in one and two noise steps, then calibrates an additional spike using the ordinary return probability of the target sphere. A binary type estimate makes the spike affordable at the critical stationary moment. We give a localized norm inequality showing that the critical moment of the normalized weighting density diverges exponentially in the same family. These conclusions concern finite Boolean noise chains.
\end{abstract}

\section{Setting and result}
Let $\mu$ be uniform on $\{-1,1\}^L$, let $P=P_{1/2}$ keep each sign with probability $3/4$ independently, and let $h>0$ with $\E_\mu hPh=1$. Put
\[
 g=Ph,\quad \pi=hg,\quad \nu=\pi\mu,\quad U=h/g,\quad r=U/\E_\mu U,
 \qquad K(x,y)=P(x,y)h(y)/g(x).
\]
Thus $\pi$ and $r$ are distinct probability densities. The stationary edge measure $\mu(x)P(x,y)h(x)h(y)$ is symmetric. $\operatorname{gap}(K)=1-\lambda_2(K)$ uses the second largest eigenvalue. Let $m(x)=L^{-1}\sum_i x_i$. A sphere $S=\{m=d\}$ is contained in the ball of radius $(1-d)L/2$ about the all-plus vertex. All function norms below use $\mu$ unless a different measure is indicated.

\begin{theorem}\label{main}
There exist fixed $d,\beta,c>0$ with $0<d<1$ and an infinite arithmetic progression of dimensions such that the following holds. For every fixed $b<1$, $b>0$, and every sufficiently large $L$ on the progression there is a positive radial potential satisfying
\[
 \pi\ge b,\qquad \norm{\pi-1}_3\le e^{-cL},\qquad
 \E_\mu\pi\chi_T=0\quad (1\le |T|\le\lfloor\beta L\rfloor).
\]
Here $\chi_T(x)=\prod_{i\in T}x_i$. In particular $\E_\mu\pi^3\le2$ eventually. The sphere $S=\{m=d\}$ has $\nu(S)\le e^{-cL}$ and
\[
 \sup_{x\in S}K(x,S^c)\le e^{-cL},\qquad
 \operatorname{gap}(K)\le e^{-cL},\qquad \norm{r}_3\ge e^{cL}.
\]
Exponents may be decreased to use one common $c$.
\end{theorem}

The marginal condition exceeds any fixed $C\sqrt L$ for large $L$. This refutes a dimension-uniform positive gap and a dimension-uniform critical weighting moment under the stationary assumptions of \cite{mom}. The sufficient escape estimate for radii up to $L/96$ at half noise in \cite{local} coexists with the trapping example at a larger fixed radius fraction. In particular that sufficient estimate cannot be extended to every fraction below $1/2$. We do not locate the optimal radius.

\section{Cost of returning to the same sphere}
For $|s|<1$ put
\[
 I(s)=\tfrac12[(1+s)\log(1+s)+(1-s)\log(1-s)],
 \qquad J(d)=\tfrac34 I(4d/3)-I(d).
\]
Fix a small positive rational $d$. Take $L$ along common multiples of the denominators of $3/8\pm d/2$, $1/8$, and $(1+d)/2$. Write $q=P(x,S)$ for $x\in S$, which is constant by permutation symmetry.
\begin{lemma}\label{type}
$\mu(S)\le e^{-LI(d)}$ and, with constants depending on $d$,
\[
 e^{-LJ(d)-O(\log L)}\le q\le e^{-Ld^2/8}.
\]
Moreover $J(d)=d^2/6+O(d^4)$ and $I(d)=d^2/2+O(d^4)$ as $d\to0$.
\end{lemma}
\begin{proof}
The first bound is the binomial Chernoff estimate, also obtained from $\binom{L}{j}\le e^{LH(j/L)}$ with binary entropy $H$. Under an ordinary stationary noise edge the sign pair probabilities, in order $++, +-, -+, --$, are $(3/8,1/8,1/8,3/8)$. Prescribe the joint empirical frequencies $(3/8+d/2,1/8,1/8,3/8-d/2)$. Both endpoints then lie in $S$. The multinomial formula and $\log n!=n\log n-n+O(\log(n+2))$ bound the probability of this joint type below by $\exp[-(3L/4)I(4d/3)-O(\log L)]$. The probability of all edges within $S$ is $\mu(S)q$, giving the lower bound. For independent signs, possibly of different means, the one-sided deviation probability of their average by $a$ is at most $e^{-La^2/2}$ \cite{hoeff}. Here the mean from the sphere is $d/2$. Finally $I(s)=s^2/2+O(s^4)$ by Taylor expansion at zero.
\end{proof}

\section{A depression, two compensation bands and a spike}
Set $w=d/100$ and $k=d^3$. Choose and fix the rational $d$ sufficiently small that
\[
 3k<w^2/8,\quad 2k<d^2/500,\quad d^2/8>2k,\quad
 J(d)-I(d)/3<k/2,\quad J(d)+k<2I(d)/3 .
\]
All these conditions follow from Lemma \ref{type} by taking $d$ smaller. Small exponents in this construction are fixed before $L$ increases. Define disjoint sets
\[
 B=\{|m-d/2|\le w\},\quad C=\{|m-d/4|\le w\},
 \quad D=\{|m-d|\le4w\}.
\]
The background potential $H_0$ equals $e^{-kL}$ on $B$, $e^{kL}$ on $C\cup D$ and one elsewhere. With $S\subset D$, define
\[
 M=e^{-kL/2}/q,\qquad H=H_0+M\one_S,\qquad W=H\,PH.
\]
By the upper bound on $q$, $M\ge e^{kL}$ eventually.

\paragraph{Uniform floor.}
The mean magnetization after ordinary noise from $x$ is $m(x)/2$. For $x\in B$, the mean lies within $w/2$ of $d/4$, so $P(x,C)\ge1-2e^{-Lw^2/8}$. The factor $e^{kL}$ on $C$ compensates $e^{-kL}$ on $B$. For $x\in D$, $H_0(x)\ge e^{kL}$ and $PH_0(x)\ge e^{-kL}$. Outside $B\cup D$ the mean is at distance greater than $2w$ from $d/2$, and $P(x,B)\le2e^{-Lw^2/2}$. There $H_0\ge1$ and $PH_0\ge1-P(x,B)$. Consequently $\inf W\ge1-O(e^{-Lw^2/8})$ by positivity. The two compensation bands ensure this bound also at vertices adjacent to the modified region.

\paragraph{Retention.}
For $x\in S$ the landing mean is exactly $d/2$, so
$PH_0(x)\le e^{-kL}+2 e^{kL-Lw^2/2}\le2e^{-kL}$ eventually.
The added spike contributes $Mq=e^{-kL/2}$ to $PH(x)$ entirely inside $S$. It follows row by row that
\[
 K_H(x,S^c)\le PH_0(x)/(Mq)\le2e^{-kL/2}\qquad(x\in S).
\]

\paragraph{Critical stationary cost.}
Write $E=B\cup C\cup D$. Every member has $m(x)\ge d/8$, so $\mu(E)\le e^{-Ld^2/128}$. By contraction of $P$ on $L^3$ and $\|H_0\|_\infty\le e^{kL}$,
\[
 \|H_0PH_0-1\|_3\le (1+e^{kL})e^{kL}\mu(E)^{1/3}.
\]
This decays exponentially by the choice of $k$. Outside $S$, the additional term is exactly $H_0 M P\one_S$. Ordinary Boolean hypercontractivity at correlation $1/2$ gives $\|P\one_S\|_3\le\|\one_S\|_{3/2}=\mu(S)^{2/3}$ \cite{bonami}. Thus its norm is at most $\exp[L(k/2+J(d)-2I(d)/3)+O(\log L)]$. On $S$, $PH\le3e^{-kL/2}$ and $H\le 2M$ eventually. In particular
\[
 \|\one_S W\|_3\le6 e^{-kL}\mu(S)^{1/3}/q
 \le 6\exp[L(J(d)-I(d)/3-k)+O(\log L)] .
\]
This too decays exponentially. Since $\|\one_S\|_3=\mu(S)^{1/3}$, the three bounds imply $\|W-1\|_3=O(e^{-c_0L})$ for some $c_0>0$. For $Z=\E W$ it follows that $Z=1+O(e^{-c_0L})$. Now $h_0=H/\sqrt Z$ gives $\pi_0=W/Z$ with the same convergence and $\inf\pi_0\ge1-o(1)$. Normalization does not change the transition law.

\section{Exact matching and the spectral gap}
For fixed $\beta\in(0,1/2)$ with binary entropy $H(\beta)<c_0/2$ in natural units, subtract the nonconstant Walsh projection from $\pi_0$:
\[
 v=\pi_0-\sum_{1\le|T|\le\beta L}(\E\pi_0\chi_T)\chi_T.
\]
There are at most $e^{LH(\beta)}$ summands by the binomial bound and each coefficient is at most $\|\pi_0-1\|_1$. Thus $\|v-\pi_0\|_\infty=O(e^{-c_0L/2})$. The density $v$ is radial with mean one; it satisfies the stated Fourier identities, floor and cubic convergence.

We include the stability argument of \cite[Lemma 10.4]{mom} to check that exact matching preserves the trap even if the gap is small. For a positive target $u$, minimize
\[
 \tfrac12\sum_{x,y}\mu(x)P(x,y)e^{z_x+z_y}-\sum_x\mu(x)u(x)z_x .
\]
The Hessian is positive definite by the diagonal terms. If the maximum coordinate tends to infinity, its diagonal term dominates the linear part. When the maximum stays bounded above and a coordinate tends to minus infinity, its positive linear weight forces divergence. Thus there is a unique minimizer, solving $e^z P e^z=u$. Positive definiteness gives smooth dependence. Along $u_t=(1-t)\pi_0+t v$, write $h_t=e^{z_t}$ and $K_t$ for its Doob kernel. Logarithmic differentiation gives $(I+K_t)\dot z_t=(v-\pi_0)/u_t$. For large $L$, $u_t\ge b_*=(1+b)/2$. The stationary law has smallest atom at least $2^{-L} b_*$. The spectrum of $K_t$ is in $[0,1]$ by positive semidefiniteness of Boolean noise and reversible conjugation. For $T=(I-K_t)/2$ the absolute row sums are at most one and its stationary $L^2$ norm is at most $1/2$. Hence
\[
 \|T^j\|_{\infty\to\infty}\le\min(1,2^{(L/2)-j}/\sqrt{b_*}),
 \qquad (I+K_t)^{-1}=\tfrac12\sum_{j\ge0}T^j .
\]
The inverse norm is $O_b(L+1)$ independent of the gap. Integration shows $\|\log(h_1/h_0)\|_\infty=O_b((L+1)e^{-c_0 L/2})$. Uniqueness gives radiality. Every transition changes by a factor between $e^{-2\eta}$ and $e^{2\eta}$ for this logarithmic error $\eta$. In particular the exit bound persists, and $h=h_1$ has stationary density $v$.

Since $\nu(S)\le\|v\|_3\mu(S)^{2/3}$, its mass is exponentially small. The centered indicator of $S$ has Dirichlet energy at most $\nu(S)\sup_{S}K(x,S^c)$ and variance $\nu(S)(1-\nu(S))$. This proves the gap upper bound. Starting anywhere on $S$, survival for $n$ steps is at least $(1-\sup_S K(x,S^c))^n$. In particular the mean exit time grows exponentially.

\section{The critical weighting moment}
We prove the localized estimate needed to distinguish a large ratio at one vertex from divergence of its critical moment. For each fixed $\alpha<1/2$, $\alpha\ge0$, there exists $a_\alpha>0$ such that on every subset $A$ of a ball of radius $\alpha L$,
\[
 \|K_{AA}\|_{L^2(\nu)\to L^2(\nu)}
 \le e^{-a_\alpha L}\|\one_A U\|_3.                                      \tag{1}
\]
In one coordinate use norms of uniform probability, $p=3/2$, $p'=3$. Encode minus signs by $j=1$, plus by $j=0$. Define $D_t(j)=e^{t(1/2-j)}$ as a multiplication operator and $N(t)=\|D_t P D_t\|_{p\to p'}$. Positivity and Cauchy--Schwarz for its positive semidefinite quadratic form identify this with the maximal quadratic form over nonnegative unit $p$-vectors. At zero the unique nonnegative maximizer is one. Indeed, normalize the mean of a nonzero pair to one and write its entries as $1\pm x$, $0\le x\le1$. Its noise energy is $1+x^2/2$. The derivative of $[(1+x)^{3/2}+(1-x)^{3/2}]/2$ is at least $3x/4$ since $\sqrt{1+x}+\sqrt{1-x}\le2$. This exceeds the derivative of $(1+x^2/2)^{3/4}$ for $x>0$; at zero both functions are one.

The derivative of $\ip{1}{D_t P D_t1}$ at zero vanishes by centering. Nonnegative maximizers for $t\downarrow0$ approach one by compactness and uniqueness. Taylor expansion in $t$ is uniform on the unit sphere, so comparison with the value at zero and testing on one yield $\log N(t)=o(t)$. The $p$-to-$p'$ tensor norm is at most $N(t)^L$, by successive coordinate estimates and Minkowski's mixed-norm inequality ($p\le p'$). On the centered ball the inverse product multiplier is at most $e^{-t(1/2-\alpha)L}$; another center follows by translation. Applying the product bound to both copies of a supported function gives
\[
 \ip{f}{P f}\le[N(t)e^{-2t(1/2-\alpha)}]^L\|f\|_{3/2}^2 .
\]
Choose small $t>0$ so the bracket is below one. Substitute $f=\sqrt U\,z$, use $\|f\|_{3/2}\le\|\one_A U\|_3^{1/2}\|z\|_2$ for supported $z$, and conjugate $K_{AA}$ by $\sqrt{\pi}$ to obtain (1).

Apply (1) with $\alpha=(1-d)/2$ and $A=S$. The Rayleigh quotient of $\one_S$ is at least $1-\sup_S K(x,S^c)$, so $\|\one_S U\|_3$ grows exponentially. Normalizing $U$ by its mean preserves this conclusion. Indeed $\pi=hg\ge b$ gives $g\ge b P(1/g)\ge b/(Pg)$ by Jensen. With $Q=P_{\sqrt{1/2}}$ we have $\|Qh\|_2=1$ and $Pg=P_{(1/2)^{3/2}}Qh$. Hypercontractivity gives $\|1/g\|_9\le1/b$. Consequently $\E U=\E(\pi/g^2)\le\|\pi\|_3\|1/g\|_3^2\le 2^{1/3}/b^2$. This completes the proof of Theorem \ref{main}.

\section{Relation to earlier results and review context}
Bonami's inequality supplies ordinary-cube norm estimates. Polyanskiy and Samorodnitsky developed improved hypercontractivity on small supports \cite{ps}. Here the stationary density and the transition row both change with the potential. The earlier construction in \cite[Section 11]{mom} has an exponentially small noise denominator at a vertex but preserves the ordinary spectrum asymptotically. The collective self-return of an entire sphere, calibrated by its conditional type probability, changes the conclusion. The shell widths and exponents here are sufficient rather than optimized. The construction requires arbitrarily large dimensions, gives a subsequence at fixed correlation, and supplies no continuum or physical spectral statement. The literature comparison does not establish global priority.

\paragraph{Model assessments and review context.} OpenAI Codex assisted with the construction, proof and manuscript. The analytic audit shares the preparation context. No proof assistant or numerical oracle was used. The first exact-PDF assessment, its objections, identity information and score belong to the version-specific record. Earlier scores concern different artifacts.

\begin{thebibliography}{9}
\bibitem{hoeff} W. Hoeffding, Probability inequalities for sums of bounded random variables, J. Amer. Statist. Assoc. 58 (1963), 13--30. \url{https://doi.org/10.1080/01621459.1963.10500830}.
\bibitem{bonami} A. Bonami, Etude des coefficients de Fourier des fonctions de $L^p(G)$, Ann. Inst. Fourier 20 (1970), 335--402. \url{https://doi.org/10.5802/aif.357}.
\bibitem{ps} Y. Polyanskiy and A. Samorodnitsky, Improved log-Sobolev inequalities, hypercontractivity and uncertainty principle on the hypercube, J. Funct. Anal. 277 (2019). \url{https://arxiv.org/abs/1606.07491}.
\bibitem{mom} L. Eriksson, Moment transfer and spectral bounds for Boolean Doob kernels, AIRR working paper, v3 (2026). \url{https://airr.science/papers/ARR-2026-4T35VQ6MZB8F99F4/versions/v3/}.
\bibitem{local} L. Eriksson, Local retention and spherical reconstruction for Boolean Doob chains, AIRR working paper, v1 (2026). \url{https://airr.science/papers/ARR-2026-6E2HS34Q2Z8349BJ/versions/v1/}.
\end{thebibliography}
\end{document}
